\section{Kết luận}\label{sec:conclusion}

Gom các chương trình sai theo những test mà chúng trượt là một cách hấp dẫn để đưa phản hồi vĩ mô cho giảng viên, và các luật NẾU--THÌ của nó đọc lên giống lời khuyên dạy học. Khi đánh giá bằng bản sửa tối thiểu đã kiểm chứng từ chính chương trình của sinh viên và bằng mutant một cơ chế có kiểm soát, OAV theo test case hóa ra trộn lẫn cơ chế một cách hệ thống: nó nhận ra các lỗi mức output chiếm đa số và gộp phần lớn phần còn lại (RQ1). Bằng chứng không cần nhãn về cách output lệch khỏi đáp án khôi phục được nhiều hơn, rõ nhất khi các cơ chế cân bằng (RQ2), và nó hỗ trợ những luật chuyển được giữa các bài tập (RQ3); ngược lại, embedding mã đa dụng gom chương trình theo nguồn gốc (RQ4). Vì mọi nhãn trong các benchmark đều suy ra từ bằng chứng thực thi và có thể suy lại, các benchmark này dùng lại được để đánh giá biểu diễn mới, bộ chẩn đoán dựa trên LLM hoặc test chủ động, và công cụ cho giảng viên cho thấy các kết quả chuyển thành phản hồi biết từ chối trả lời khi bằng chứng yếu như thế nào.

Ba bước sẽ củng cố các kết luận này: một lần audit nhãn độc lập bởi hai người chấm, lặp lại trên khóa học, ngôn ngữ và chính sách chấm khác, và một nghiên cứu trong lớp học xem giả thuyết cơ chế có thay đổi điều giảng viên dạy lại hay không. Test thích ứng để phá các va chạm chữ ký còn lại trước khi giả thuyết đến tay giảng viên là hướng phương pháp tự nhiên tiếp theo. Khẳng định về niềm tin của sinh viên cần nghiên cứu với sinh viên; các công cụ được công bố ở đây nhằm làm cho việc thiết kế những nghiên cứu đó rẻ hơn.
